{\color{magenta}
\chapter{Relevant Parameters for the Polomac}
\label{chap:parameters}

Three dimensionless (or characteristic) parameters govern the MHD
dynamics of the Polomac.  For each we first introduce the underlying
physics, then evaluate the parameter using design values
from~\cite{Elio2024}, and finally state which assumption of subsection \ref{subsec:mhd_assumptions} it justifies.

Throughout this chapter we use the following Polomac
values~\cite{Elio2024}: hydrogen plasma ($m_i=1.67\times10^{-27}$~kg,
$Z=1$), $n=10^{20}$~m$^{-3}$, $T_i\approx T_e=100$~eV,
$B=0.2$--$0.3$~T, $L\sim0.04$--$0.3$~m.

% ------------------------------------------------------------
\section{Alfv\'en Velocity and Waves}
\label{sec:alfven}

\subsection*{Physics}

Alfv\'en waves are transverse, incompressible oscillations that
propagate along magnetic field lines, driven by magnetic tension and
resisted by plasma inertia.  To derive their dispersion relation we
linearize the ideal MHD equations around a static uniform equilibrium
$\vec{B}=B_0\hat{e}_z$, $\varrho=\varrho_0$, $\vec{V}=0$,
$P=P_0$, writing each quantity as its equilibrium value plus a small
perturbation (e.g.\ $\vec{B}=B_0\hat{e}_z+\vec{B}_1$).  To first
order, the momentum equation and induction equation (Ohm's law combined
with Faraday's law) give
\begin{align}
    \varrho_0\frac{\partial\vec{V}_1}{\partial t}
        &= \frac{1}{\mu_0}(\nabla\wedge\vec{B}_1)\wedge\vec{B}_0,
    \label{eq:alfven_mom}\\
    \frac{\partial\vec{B}_1}{\partial t}
        &= \nabla\wedge(\vec{V}_1\wedge\vec{B}_0).
    \label{eq:alfven_ind}
\end{align}
Assuming plane-wave solutions $\sim e^{i(\vec{k}\cdot\vec{x}-\omega t)}$
with $\vec{k}=k\hat{e}_z$ and perturbations perpendicular to
$\vec{B}_0$, one obtains the dispersion relation
\begin{equation}
    \omega^2 = k_\parallel^2\, v_A^2,
    \label{eq:alfven_disp}
\end{equation}
where $k_\parallel=\vec{k}\cdot\hat{b}$ is the wavenumber along the
field and
\begin{equation}
    \boxed{v_A = \frac{B}{\sqrt{\mu_0\varrho}}}
    \label{eq:va}
\end{equation}
is the \emph{Alfv\'en velocity}.  The waves are non-dispersive and
propagate without compression.

\subsection*{Evaluation for the Polomac}

With $B=0.25$~T (midpoint of operating range) and
$\varrho = nm_i = 10^{20}\times1.67\times10^{-27}
= 1.67\times10^{-7}$~kg\,m$^{-3}$:
\[
    v_A
    = \frac{0.25}{\sqrt{4\pi\times10^{-7}
                        \times1.67\times10^{-7}}}
    = \frac{0.25}{4.58\times10^{-7}}
    \approx 5.5\times10^5\,\text{m\,s}^{-1}.
\]
The corresponding MHD timescale for $L=0.3$~m is
$\tau_A = L/v_A \approx 5.5\times10^{-7}$~s.

\subsection*{Significance and justification of MHD assumptions}

\paragraph{Assumption~4: neglect displacement current.}
\[
    \frac{v_A}{c}
    = \frac{5.5\times10^5}{3\times10^8}
    \approx 1.8\times10^{-3} \ll 1. \quad
\]
MHD dynamics are non-relativistic; the displacement current in
Amp\`ere's law is negligible.

\paragraph{Assumption~1: electron inertia negligible.}
The ion thermal velocity is
\[
    V_{Ti}
    = \sqrt{\frac{2\times100\,\text{eV}\times1.6\times10^{-19}\,
                  \text{J\,eV}^{-1}}
                 {1.67\times10^{-27}\,\text{kg}}}
    = \sqrt{1.916\times10^{10}}
    \approx 1.38\times10^5\,\text{m\,s}^{-1}.
\]
The ECR heating frequency implies $\omega_{ce}=2\pi\times4\times10^9
=2.51\times10^{10}$~rad\,s$^{-1}$.  Using $L=0.3$~m:
\[
    \frac{\omega_{MHD}}{\omega_{ce}}
    = \frac{V_{Ti}/L}{\omega_{ce}}
    = \frac{1.38\times10^5/0.3}{2.51\times10^{10}}
    = \frac{4.6\times10^5}{2.51\times10^{10}}
    \approx 1.8\times10^{-5} \ll 1. \quad
\]
Ion dynamics are slow compared to the electron cyclotron motion;
electron inertia can be neglected on MHD timescales.

% ------------------------------------------------------------
\section{Lundquist Number}
\label{sec:lundquist}

\subsection*{Physics}

The Lundquist number compares the resistive diffusion timescale
$\tau_R = \mu_0 L^2/\eta$ to the Alfv\'en timescale $\tau_A=L/v_A$:
\begin{equation}
    \boxed{S \coloneqq \frac{\tau_R}{\tau_A} = \frac{\mu_0 L v_A}{\eta}.}
    \label{eq:lundquist}
\end{equation}
Equivalently, $S$ is the magnetic Reynolds number evaluated at the
Alfv\'en speed.  When $S\gg1$, resistive diffusion is slow on MHD
timescales and the magnetic flux is effectively frozen into the
fluid — the ideal MHD limit.

\subsection*{Evaluation for the Polomac}

Using $\eta\sim10^{-6}$~$\Omega$m (Section~\ref{sec:mhd_assumptions}),
$v_A=5.5\times10^5$~m\,s$^{-1}$, and $L=0.3$~m:
\[
    S = \frac{4\pi\times10^{-7}\times0.3\times5.5\times10^5}
             {10^{-6}}
    \approx 2\times10^5.
\]

\subsection*{Significance and justification of MHD assumptions}

\paragraph{Assumption~3: zero resistivity.}
$S\approx2\times10^5\gg1$ confirms that resistive diffusion is
negligible on MHD timescales, justifying $\eta\to0$.  This value is
modest compared to large tokamaks ($S\sim10^6$--$10^8$), reflecting the
prototype's small size and low temperature.  Since $\eta\propto
T^{-3/2}$, a reactor-scale Polomac at keV temperatures would reach
$S\sim10^6$.

For the static equilibrium analysis of Chapter~\ref{chap:equilibrium}
($\vec{V}=0$), resistive effects are absent regardless, so even a
moderate $S$ is sufficient to justify the framework.

% ------------------------------------------------------------
\section{Plasma \texorpdfstring{$\beta$}{Beta}}
\label{sec:beta}

\subsection*{Physics}

The plasma $\beta$ is the ratio of kinetic plasma pressure to magnetic
pressure:
\begin{equation}
    \boxed{\beta \coloneqq \frac{p}{B^2/(2\mu_0)}.}
    \label{eq:beta}
\end{equation}
It measures confinement efficiency: a higher $\beta$ means more fusion
power per unit of magnetic field energy invested.  Conventional
tokamaks are typically limited to $\beta lesssim 0.1$ by current-driven
instabilities.

\subsection*{Evaluation for the Polomac}

From~\cite{Elio2024} (Table~1), plasma pressure is $1.6$--$16$~kPa.
The magnetic pressure at the operating field range is:
\[
    \frac{B^2}{2\mu_0}\bigg|_{B=0.2\,\text{T}}
    = \frac{0.04}{2\times1.257\times10^{-6}}
    \approx 15.9\,\text{kPa},
\]
\[
    \frac{B^2}{2\mu_0}\bigg|_{B=0.3\,\text{T}}
    = \frac{0.09}{2\times1.257\times10^{-6}}
    \approx 35.8\,\text{kPa}.
\]
This gives:
\[
    \beta \in [0.04,\;1].
\]

\subsection*{Significance for the Polomac}

The predominantly poloidal field topology is claimed to support higher
$\beta$ than toroidal devices because the closed poloidal field lines
reduce current-driven instability drive~\cite{Elio2024}.  If the upper
range $\beta\to1$ is achievable, the Polomac could confine fusion-grade
pressures with a comparatively modest magnetic field, reducing both
engineering complexity and power consumption.  These are design targets
and have not yet been experimentally demonstrated.
}